\documentclass[12pt]{article}
\usepackage[margin=1in]{geometry}

\title{Design Specification\\FlightSound}
\author{Group 05}
\date{May 11, 2026}

\begin{document}

\maketitle
\newpage

\tableofcontents
\newpage

\section*{Version History}
\addcontentsline{toc}{section}{Version History}

\begin{tabular}{|l|l|l|l|}
\hline
Version & Date & Author & Description \\
\hline
1.0 & May 11, 2026 & Group 06 & Initial Draft \\
\hline
\end{tabular}

\newpage

\section{Page and Component Breakdown}

This section describes every page, panel, and tool that makes up the FlightSound application.

\subsection{Main Map Page}
The main map page is the primary screen of the application. It loads on startup and serves as the central interface for all user interactions.

\begin{itemize}
\item Displays a full-screen interactive map rendered by the Google Maps JavaScript API.
\item Defaults to LAX as the starting airport when the page first loads.
\item Shows aircraft markers as icons positioned at their current coordinates on the map.
\item Shows color-coded noise circles around each aircraft: green for low noise, yellow for medium, and red for high.
\item Contains the navigation bar, search bar, and all toggle controls.
\item Sends HTTP requests to the Java backend at port 8080 to retrieve flight and weather data.
\end{itemize}

\subsection{Airport Selection}
The airport selection tool is embedded in the navigation bar at the top of the main page.

\begin{itemize}
\item Allows users to type an airport name or IATA code into a search input.
\item Sends the selected airport code to the backend API to fetch nearby flight data.
\item Recenters the map on the selected airport after a selection is made.
\item Validates the input to only accept known IATA airport codes.
\end{itemize}

\subsection{Aircraft Markers}
Aircraft markers are rendered on the map for each active flight returned by the backend.

\begin{itemize}
\item Each marker is positioned at the flight's current latitude and longitude.
\item Clicking a marker opens the aircraft details side panel.
\item Markers are accompanied by noise zone circles that visually represent the estimated noise impact area.
\item Marker positions are updated when new flight data is fetched.
\end{itemize}

\subsection{Aircraft Details Side Panel}
The side panel appears when a user clicks on an aircraft marker or a row in the results table.

\begin{itemize}
\item Displays the flight number, airline name, aircraft type, altitude, and speed.
\item Shows the estimated noise level in decibels and a descriptive label such as Whisper, Conversation, or Jet Engine.
\item Shows the departure and arrival airport codes.
\item Displays the noise zone classification: low, medium, or high.
\end{itemize}

\subsection{Noise Heatmap and Zone Circles}
The noise visualization tools allow users to see estimated noise impact areas on the map.

\begin{itemize}
\item Color-coded circles are drawn around each aircraft marker representing the estimated noise impact zone.
\item Low noise zones are shown in green, medium in yellow, and high in red.
\item The circles can be toggled on or off using a control button on the map.
\item Noise levels are calculated by the backend using a 3D slant-distance model combining horizontal distance and altitude.
\end{itemize}

\subsection{Weather Panel}
The weather panel displays current conditions for the selected airport.

\begin{itemize}
\item Retrieves weather data from the Open-Meteo API via the backend weather endpoint.
\item Displays temperature, wind speed, wind direction, humidity, and general weather condition.
\item Weather data is fetched using the airport's latitude and longitude coordinates.
\item Displayed in the side panel below aircraft details or as a standalone view.
\end{itemize}

\subsection{Google Sign-In and Authentication}
FlightSound includes a Google OAuth sign-in system managed by the AuthManager module in the frontend.

\begin{itemize}
\item Users can sign in using their Google account via the Google Identity Services API.
\item On sign-in, the frontend decodes the Google JWT token to extract the user's Google ID, email, and display name.
\item The user's session is stored in the browser's localStorage under the key fs\_user.
\item On first sign-in, a new user record is created in the SQLite database via the backend auth endpoint.
\item Returning users are looked up by their Google ID and their session is restored.
\end{itemize}

\subsection{User Preferences}
The preferences panel allows signed-in users to save personal settings.

\begin{itemize}
\item Users can save a preferred airport so the map defaults to that location on load.
\item Users can save a preferred address which is geocoded and used as their noise reference point.
\item Preferences are stored in the SQLite database and retrieved on sign-in via the preferences endpoint.
\item The noise calculation mode switches between airport-based, user-location-based, and address-based depending on the user's preference.
\end{itemize}

\subsection{Help Page}
The Help page is accessible from the navigation bar.

\begin{itemize}
\item Explains how noise levels are estimated and what each color zone means.
\item Describes how to use the airport search and map controls.
\item Includes a brief FAQ section covering common questions about the data and its limitations.
\item Clarifies that noise estimates are approximations and not certified acoustic measurements.
\end{itemize}

\subsection{Backend API Endpoints}
The Java backend runs on port 8080 and exposes the following HTTP endpoints:

\begin{tabular}{|l|l|p{6cm}|}
\hline
Endpoint & Method & Description \\
\hline
/api/noise?airport=\{iata\} & GET & Returns flight data and estimated noise levels for the given airport. \\
\hline
/api/weather?airport=\{iata\} & GET & Returns current weather for the given airport using Open-Meteo. \\
\hline
/api/weather?lat=\{lat\}\&lng=\{lng\} & GET & Returns weather for arbitrary coordinates. \\
\hline
/api/auth & POST & Handles Google Sign-In, creates or retrieves a user record. \\
\hline
/api/users & GET/PUT & Retrieves or updates user account information. \\
\hline
/api/user/preferences & GET/PUT & Retrieves or updates saved preferences for a signed-in user. \\
\hline
\end{tabular}

\subsection{SQLite Database}
FlightSound uses a local SQLite database file (users.db) to store user account and preferences data.

\begin{tabular}{|l|p{4cm}|p{6cm}|}
\hline
Table & Key Fields & Description \\
\hline
users & google\_id (PK), email, username, created\_at & Stores registered user accounts linked to Google Sign-In. \\
\hline
preferences & google\_id (FK), preferred\_airport, preferred\_address, address\_lat, address\_lng & Stores saved user preferences for airport and noise reference location. \\
\hline
\end{tabular}

\newpage

\section{Docker Dependencies}

This section lists all packages, libraries, and dependencies required to run FlightSound inside a Docker container.

\subsection{Backend Container}

The backend runs as a Java application. The following are required:

\begin{tabular}{|l|l|p{6cm}|}
\hline
Dependency & Version & Purpose \\
\hline
Java JDK & 17 or higher & Runtime and compilation environment for the Java backend. \\
\hline
Google Gson & 2.10.1 & JSON serialization and deserialization for API responses and request parsing. \\
\hline
sqlite-jdbc & 3.45.1.0 & JDBC driver that allows the Java backend to connect to and query the SQLite database file. \\
\hline
com.sun.net.httpserver & Built into JDK & Lightweight HTTP server used to expose the backend API endpoints. No additional download required. \\
\hline
java.net.http.HttpClient & Built into JDK 11+ & Used to send HTTP requests to the AirLabs API and Open-Meteo API. No additional download required. \\
\hline
\end{tabular}

Environment variables that must be set in the container:

\begin{tabular}{|l|p{9cm}|}
\hline
Variable & Description \\
\hline
FLIGHTLABS\_API\_KEY & API key for the AirLabs real-time flight data API. \\
\hline
GOOGLE\_MAPS\_API\_KEY & API key for the Google Maps JavaScript API used by the frontend. \\
\hline
\end{tabular}

\subsection{Frontend Container}

The frontend is a set of static files (HTML, CSS, JavaScript) that require a web server to be served. No build step or package manager is needed.

\begin{tabular}{|l|l|p{6cm}|}
\hline
Dependency & Version & Purpose \\
\hline
Nginx & Latest stable & Serves the static frontend files (HTML, CSS, JS) to the browser. \\
\hline
Bootstrap & 5.x (CDN) & Loaded via CDN link in the HTML. No local installation required inside the container. \\
\hline
Google Maps JavaScript API & Latest (CDN) & Loaded via script tag using the API key. No local installation required. \\
\hline
Google Identity Services & Latest (CDN) & Loaded via script tag for Google Sign-In. No local installation required. \\
\hline
\end{tabular}

\subsection{External APIs (No Docker Installation Required)}
The following services are called over HTTPS at runtime and do not need to be installed in the container:

\begin{tabular}{|l|p{10cm}|}
\hline
Service & Description \\
\hline
AirLabs API & Provides real-time flight data. Requires FLIGHTLABS\_API\_KEY environment variable. \\
\hline
Open-Meteo API & Provides current weather data by coordinates. No API key required. \\
\hline
Google Maps JavaScript API & Renders the interactive map. Requires GOOGLE\_MAPS\_API\_KEY. \\
\hline
Google Identity Services & Handles Google OAuth sign-in. Requires a configured Google OAuth client ID. \\
\hline
\end{tabular}

\newpage

\section{Glossary}

\begin{tabular}{|l|l|}
\hline
Acronym & Definition \\
\hline
API & Application Programming Interface \\
CDN & Content Delivery Network \\
CSS & Cascading Style Sheets \\
dB & Decibel \\
HTML & HyperText Markup Language \\
HTTP & HyperText Transfer Protocol \\
HTTPS & HyperText Transfer Protocol Secure \\
IATA & International Air Transport Association \\
JDBC & Java Database Connectivity \\
JDK & Java Development Kit \\
JSON & JavaScript Object Notation \\
JWT & JSON Web Token \\
OAuth & Open Authorization \\
SQLite & Self-Contained Lightweight SQL Database Engine \\
UI & User Interface \\
URL & Uniform Resource Locator \\
\hline
\end{tabular}

\end{document}
